\chapter{Dilation and Curettage (D\&C)}

\section*{Overview}
Dilation and curettage (D\&C) involves dilating the cervix and scraping the uterine lining (endometrium) with a curette to remove tissue.

\section*{Alternative Names}
D\&C, Uterine scrapings

\section*{Principle}
The cervix is dilated and a curette is used to scrape the uterine cavity, removing the endometrial tissue for diagnosis or to treat retained products of conception, abnormal bleeding, or other conditions.
\section*{History}
Developed in the 19th century as a diagnostic and therapeutic tool for uterine bleeding and managing incomplete miscarriages.

\section*{Why It Is Done}
Ordered to diagnose endometrial hyperplasia or cancer, treat heavy uterine bleeding, or clear retained tissue after a miscarriage or delivery.

\section*{Is This a Primary or Secondary Test or Procedure}
Primary therapeutic method for heavy uterine bleeding refractory to medical therapy and for retained products of conception.

\section*{How It Is Performed}
Performed under sedation or regional anesthesia. The cervix is dilated using metal dilators, and a curette (or suction device) is inserted into the uterus to scrape or suction the endometrial lining.

\section*{Preparation}
Fast for 8 hours if sedation is planned. Have a negative pregnancy test (unless clearing retained products).

\section*{Contraindications and Precautions}
Active pelvic inflammatory disease (PID), or suspected viable pregnancy (unless therapeutic abortion is intended).

\section*{Risks}
Uterine perforation, cervical damage, post-procedure pelvic infection, or Asherman's syndrome (intrauterine adhesions).

\section*{Aftercare and Recovery}
Mild cramping and spotting are common for a few days. Avoid tampons and sexual intercourse for 2 weeks to prevent infection.

\section*{Results and Interpretation}
Tissue removed is sent to pathology to inspect for abnormal endometrial cells, polyps, or retained products of conception.

\section*{Expected Results}
Successful evacuation of the uterine cavity; pathology shows benign secretory/proliferative endometrium or expected gestational tissue.

\section*{Follow-up}
Pathology review; clinical follow-up in 2 weeks.

\section*{References}
\begin{enumerate}
  \item MedlinePlus. Dilation and Curettage (D\&C). U.S. National Library of Medicine; 2025.
  \item Current Medical Diagnosis and Treatment. McGraw-Hill; 2025.
\end{enumerate}

\clearpage
